\section{ĐÁP ÁN: ĐỀ THI MẪU}

\subsection{Đưa về dạng Chuẩn tắc}
\textbf{a. Chuyển bài toán (P) về dạng chuẩn và xác định các ma trận} 

Vì tất cả các ràng buộc chính đều mang dấu "$\le$", ta cộng thêm các biến phụ không âm $x_4, x_5, x_6 \ge 0$ vào vế trái để tạo thành các đẳng thức. Ta có mô hình bài toán dạng chuẩn như sau:
\begin{alignat*}{4}
    Z = \quad & 11x_1 & {}+{} & 9x_2 & {}+{} & 15x_3 & {}+{} 0x_4 + 0x_5 + 0x_6 \rightarrow \max \\
    \text{s.t.} \quad & 5x_1 & {}+{} & 4x_2 & {}+{} & 4x_3 & {}+{} x_4 = 21000 \\
    & 3x_1 & {}+{} & 6x_2 & {}+{} & 3x_3 & {}+{} x_5 = 15000 \\
    & 3x_1 & {}+{} & 7x_2 & {}+{} & 5x_3 & {}+{} x_6 = 26000 \\
    & \rlap{$x_j \ge 0, \quad j = \overline{1,6}$}
\end{alignat*}

Từ mô hình dạng chuẩn trên, ta xác định được chi tiết các ma trận và vector cấu thành bài toán:
\begin{itemize}
    \item \textbf{Vector hệ số hàm mục tiêu:} $C = (11, 9, 15, 0, 0, 0)$.
    \item \textbf{Vector vế phải (giới hạn nguồn lực):} $b = (21000, 15000, 26000)^T$.
    \item \textbf{Ma trận điều kiện $A$ (kích thước $3 \times 6$):} Chứa các hệ số của hệ ràng buộc chính.
    \begin{equation*}
        A = \begin{pmatrix}
            5 & 4 & 4 & 1 & 0 & 0 \\
            3 & 6 & 3 & 0 & 1 & 0 \\
            3 & 7 & 5 & 0 & 0 & 1
        \end{pmatrix}
    \end{equation*}
    \item \textbf{Các vector điều kiện $A_j$:} Tương ứng với cột thứ $j$ của ma trận $A$. Ta thấy hệ vector $\{A_4, A_5, A_6\}$ tạo thành một ma trận đơn vị, do đó $x_4, x_5, x_6$ đóng vai trò là các biến cơ sở xuất phát.
\end{itemize}

\vspace{0.5cm}
\subsection{Giải bằng phương pháp đơn hình}
\textbf{b. Giải bằng phương pháp đơn hình (Trình bày từng bước)}

Để thực hiện thuật toán đơn hình, ta áp dụng các công thức nền tảng sau:
\begin{itemize}
    \item Công thức tính hệ số ước lượng: $\Delta_j = \sum_{i=1}^{m} C_i a_{ij} - c_j$.
    \item Công thức tính tỷ số đơn hình: $\lambda_i = \frac{b_i}{a_{ik}}$ (chỉ tính với các hệ số $a_{ik} > 0$).
    \item \textbf{Quy tắc Dantzig:} Chọn biến vào cơ sở là biến có $\Delta_k < 0$ và nhỏ nhất. Chọn biến ra khỏi cơ sở là biến có tỷ số $\lambda_i$ nhỏ nhất.
\end{itemize}

\textbf{Bước 1: Thiết lập Phương án cơ bản xuất phát (Bảng 1)} \\
Cho các biến không cơ sở $x_1 = x_2 = x_3 = 0$. Ta được PACB xuất phát là $x_4 = 21000$, $x_5 = 15000$, $x_6 = 26000$ với giá trị hàm mục tiêu $Z = 0$.

% BẢNG ĐƠN HÌNH 1
\begin{table}[htbp]
    \centering
    \renewcommand{\arraystretch}{1.3}
    \begin{tabular}{|c|c|c|c|c|c|c|c|c|c|}
        \hline
        \multirow{2}{*}{\textbf{Cơ sở}} & \multirow{2}{*}{$\mathbf{C_b}$} & \multirow{2}{*}{\textbf{PA}} & $x_1$ & $x_2$ & $x_3$ & $x_4$ & $x_5$ & $x_6$ & \multirow{2}{*}{\textbf{Tỷ số $\lambda$}} \\
        \cline{4-9}
        & & & 11 & 9 & 15 & 0 & 0 & 0 & \\
        \hline
        $x_4$ & 0 & 21000 & 5 & 4 & 4 & 1 & 0 & 0 & $21000/4 = 5250$ \\
        \hline
        $x_5$ & 0 & 15000 & 3 & 6 & \textbf{3} & 0 & 1 & 0 & $15000/3 = \mathbf{5000} \rightarrow$ \\
        \hline
        $x_6$ & 0 & 26000 & 3 & 7 & 5 & 0 & 0 & 1 & $26000/5 = 5200$ \\
        \hline
        \multicolumn{2}{|c|}{$\Delta_j$} & \textbf{Z = 0} & -11 & -9 & \textbf{-15} $\uparrow$ & 0 & 0 & 0 & \\
        \hline
    \end{tabular}
    \caption{Bảng đơn hình xuất phát}
\end{table}

\vspace{0.3cm}
\textbf{Giải thích chi tiết cách tính hệ số ước lượng $\Delta_j$ tại Bảng 1:}

Trong thuật toán đơn hình, hệ số ước lượng $\Delta_j$ được tính theo công thức tổng quát:
\begin{equation*}
    \Delta_j = \sum_{i=1}^{m} C_i a_{ij} - c_j
\end{equation*}

Tại Bảng đơn hình xuất phát (Bảng 1), các biến cơ sở là các biến phụ $x_4, x_5, x_6$. Do đó, cột hệ số của biến cơ sở trong hàm mục tiêu là vector không: $C_b = (0, 0, 0)^T$. 
Hệ số hàm mục tiêu tương ứng của các biến là $c = (11, 9, 15, 0, 0, 0)$.

Vì $C_b = (0, 0, 0)^T$, nên tích vô hướng $\sum C_i a_{ij}$ tại mọi cột đều bằng $0$. Từ đó, công thức tính $\Delta_j$ tại Bảng 1 được rút gọn thành: $\Delta_j = - c_j$. Chi tiết các phép tính như sau:

\begin{itemize}
    \item \textbf{Đối với các biến không cơ sở ($x_1, x_2, x_3$):}
    \begin{itemize}
        \item $\Delta_1 = (0 \times 5 + 0 \times 3 + 0 \times 3) - 11 = \mathbf{-11}$
        \item $\Delta_2 = (0 \times 4 + 0 \times 6 + 0 \times 7) - 9 = \mathbf{-9}$
        \item $\Delta_3 = (0 \times 4 + 0 \times 3 + 0 \times 5) - 15 = \mathbf{-15}$
    \end{itemize}
    
    \item \textbf{Đối với các biến cơ sở ($x_4, x_5, x_6$):} 
    Theo nguyên tắc, hệ số ước lượng của các biến đang nằm trong cơ sở luôn luôn bằng $0$. Ta có thể kiểm chứng bằng công thức:
    \begin{itemize}
        \item $\Delta_4 = (0 \times 1 + 0 \times 0 + 0 \times 0) - 0 = \mathbf{0}$
        \item $\Delta_5 = (0 \times 0 + 0 \times 1 + 0 \times 0) - 0 = \mathbf{0}$
        \item $\Delta_6 = (0 \times 0 + 0 \times 0 + 0 \times 1) - 0 = \mathbf{0}$
    \end{itemize}
\end{itemize}

Như vậy, dòng giá trị $\Delta_j$ thu được ở Bảng 1 lần lượt là: $-11, -9, -15, 0, 0, 0$.

\textit{Đánh giá Bảng 1:} 
Ta thấy $\Delta_3 = -15 < 0$ là giá trị âm nhỏ nhất, nên chọn cột $x_3$ đưa vào cơ sở. Ở cột tỷ số, $\min(5250, 5000, 5200) = 5000$ tương ứng với dòng $x_5$, nên chọn $x_5$ ra khỏi cơ sở. Phần tử trụ (chốt) là $a_{23} = 3$.

\vspace{0.3cm}
\textbf{Bước 2: Xây dựng PACB mới bằng phép biến đổi Gauss-Jordan (Bảng 2)} \\
Thực hiện biến đổi bảng theo phần tử trụ $3$:
\begin{itemize}
    \item \textbf{Dòng $x_3$ mới} = Dòng $x_5$ cũ chia cho $3$. Ví dụ: $PA = 15000/3 = 5000$.
    \item \textbf{Dòng $x_4$ mới} = Dòng $x_4$ cũ $-$ ($4 \times$ Dòng $x_3$ mới). Ví dụ: $PA = 21000 - 4(5000) = 1000$.
    \item \textbf{Dòng $x_6$ mới} = Dòng $x_6$ cũ $-$ ($5 \times$ Dòng $x_3$ mới). Ví dụ: $PA = 26000 - 5(5000) = 1000$.
\end{itemize}

% BẢNG ĐƠN HÌNH 2
\begin{table}[htbp]
    \centering
    \renewcommand{\arraystretch}{1.3}
    \begin{tabular}{|c|c|c|c|c|c|c|c|c|}
        \hline
        \multirow{2}{*}{\textbf{Cơ sở}} & \multirow{2}{*}{$\mathbf{C_b}$} & \multirow{2}{*}{\textbf{PA}} & $x_1$ & $x_2$ & $x_3$ & $x_4$ & $x_5$ & $x_6$ \\
        \cline{4-9}
        & & & 11 & 9 & 15 & 0 & 0 & 0 \\
        \hline
        $x_4$ & 0 & 1000 & 1 & -4 & 0 & 1 & -4/3 & 0 \\
        \hline
        $x_3$ & 15 & 5000 & 1 & 2 & 1 & 0 & 1/3 & 0 \\
        \hline
        $x_6$ & 0 & 1000 & -2 & -3 & 0 & 0 & -5/3 & 1 \\
        \hline
        \multicolumn{2}{|c|}{$\Delta_j$} & \textbf{Z = 75000} & 4 & 21 & 0 & 0 & 5 & 0 \\
        \hline
    \end{tabular}
    \caption{Bảng đơn hình tối ưu}
\end{table}

\textit{Kết luận:} \\
Tại Bảng 2, tất cả các hệ số ước lượng đều không âm ($\Delta_j \ge 0, \forall j$), do đó tiêu chuẩn tối ưu đã được thỏa mãn và thuật toán dừng lại. 
Phương án sản xuất tối ưu của bài toán gốc là: $x_1 = 0$, $x_2 = 0$, $x_3 = 5000$ (Nghĩa là chỉ tập trung sản xuất 5000 sản phẩm loại C). 
Tổng tiền lãi lớn nhất đạt được là $Z_{max} = 75000$ (nghìn đồng), tương đương $75$ triệu đồng.

\vspace{0.3cm}
\textit{Kết luận (2) \textbf{lấy điểm}:} \\
Vì $\Delta_j \ge 0, \forall j=\overline{1,6}$ \\
=> PATƯ của $(P_1)$ là $(0,0,5000,1000,0,1000)$ \\
=> PATƯ của $(P)$ là $(0,0,5000)$ \\
Vậy $Z = 11 \times 0 + 9 \times 0 + 15 \times 5000 = 75000$ \\

\vspace{0.3cm}
\textbf{Giải thích chi tiết các phép tính đại số trong phép biến đổi Gauss-Jordan:}

Mục tiêu của phép biến đổi Gauss-Jordan là đưa biến $x_3$ vào cơ sở thay thế cho $x_5$. Về mặt ma trận, ta cần biến đổi cột hệ số của $x_3$ (đang là vector $(4, 3, 5)^T$ ở Bảng 1) trở thành vector đơn vị $(0, 1, 0)^T$ ở Bảng 2. Phần tử trụ được chọn là $a_{23} = 3$. 

Để tiện tính toán, ta coi mỗi dòng của bảng đơn hình (bao gồm cả cột Phương Án PA) như một vector. Gọi $R_1, R_2, R_3$ lần lượt là các vector dòng $x_4, x_5, x_6$ tại Bảng 1:
\begin{align*}
    R_1^{(c\tilde{u})} &= \begin{pmatrix} 21000 & 5 & 4 & 4 & 1 & 0 & 0 \end{pmatrix} \\
    R_2^{(c\tilde{u})} &= \begin{pmatrix} 15000 & 3 & 6 & 3 & 0 & 1 & 0 \end{pmatrix} \quad \text{(Dòng chứa phần tử trụ $3$)} \\
    R_3^{(c\tilde{u})} &= \begin{pmatrix} 26000 & 3 & 7 & 5 & 0 & 0 & 1 \end{pmatrix}
\end{align*}

\vspace{0.3cm}
\textbf{Giải thích chi tiết các phép tính đại số trong phép biến đổi Gauss-Jordan:}

Mục tiêu của phép biến đổi Gauss-Jordan là đưa biến $x_3$ vào cơ sở thay thế cho $x_5$. Về mặt ma trận, ta cần biến đổi cột hệ số của $x_3$ (đang là vector $(\textcolor{blue}{4}, \textcolor{red}{3}, \textcolor{blue}{5})^T$ ở Bảng 1) trở thành vector đơn vị $(\textcolor{green!60!black}{0}, \textcolor{green!60!black}{1}, \textcolor{green!60!black}{0})^T$ ở Bảng 2. Phần tử trụ được chọn là $a_{23} = \textcolor{red}{3}$ [1].

Để tiện tính toán, ta coi mỗi dòng của bảng đơn hình (bao gồm cả cột Phương Án PA) như một vector. Gọi $R_1, R_2, R_3$ lần lượt là các vector dòng $x_4, x_5, x_6$ tại Bảng 1. Ta sẽ tô màu cột $x_3$ để làm nổi bật các thành phần cần biến đổi:
\begin{align*}
    R_1^{(c\tilde{u})} &= \begin{pmatrix} 21000 & 5 & 4 & \textcolor{blue}{4} & 1 & 0 & 0 \end{pmatrix} \quad \text{(Cần khử số \textcolor{blue}{4} thành 0)} \\
    R_2^{(c\tilde{u})} &= \begin{pmatrix} 15000 & 3 & 6 & \textbf{\textcolor{red}{3}} & 0 & 1 & 0 \end{pmatrix} \quad \text{(Dòng chứa phần tử trụ \textcolor{red}{3})} \\
    R_3^{(c\tilde{u})} &= \begin{pmatrix} 26000 & 3 & 7 & \textcolor{blue}{5} & 0 & 0 & 1 \end{pmatrix} \quad \text{(Cần khử số \textcolor{blue}{5} thành 0)}
\end{align*}

Áp dụng quy tắc xoay bảng [1], ta tính các dòng mới của Bảng 2 như sau để biến đổi cột $x_3$ thành $(0, 1, 0)^T$:

\begin{itemize}
    \item \textbf{Tính Dòng xoay mới (Dòng $x_3$):} Lấy dòng cũ chia cho phần tử trụ $\textcolor{red}{3}$ để tạo số $\textcolor{green!60!black}{1}$ ở vị trí chốt:
    \begin{equation*}
        R_2^{(m\text{ới})} = \frac{1}{\textcolor{red}{3}} R_2^{(c\tilde{u})} = \begin{pmatrix} 5000 & 1 & 2 & \textbf{\textcolor{green!60!black}{1}} & 0 & \frac{1}{3} & 0 \end{pmatrix}
    \end{equation*}
    
    \item \textbf{Tính Dòng $x_4$ mới:} Để khử số $\textcolor{blue}{4}$, ta lấy dòng $x_4$ cũ trừ đi ($\textcolor{blue}{4} \times$ Dòng $x_3$ mới):
    \begin{align*}
        R_1^{(m\text{ới})} &= R_1^{(c\tilde{u})} - \textcolor{blue}{4} \times R_2^{(m\text{ới})} \\
        &= \begin{pmatrix} 21000 & 5 & 4 & \textcolor{blue}{4} & 1 & 0 & 0 \end{pmatrix} - \begin{pmatrix} 20000 & 4 & 8 & \textcolor{blue}{4} & 0 & \frac{4}{3} & 0 \end{pmatrix} \\
        &= \begin{pmatrix} 1000 & 1 & -4 & \textbf{\textcolor{green!60!black}{0}} & 1 & -\frac{4}{3} & 0 \end{pmatrix}
    \end{align*}

    \item \textbf{Tính Dòng $x_6$ mới:} Để khử số $\textcolor{blue}{5}$, ta lấy dòng $x_6$ cũ trừ đi ($\textcolor{blue}{5} \times$ Dòng $x_3$ mới):
    \begin{align*}
        R_3^{(m\text{ới})} &= R_3^{(c\tilde{u})} - \textcolor{blue}{5} \times R_2^{(m\text{ới})} \\
        &= \begin{pmatrix} 26000 & 3 & 7 & \textcolor{blue}{5} & 0 & 0 & 1 \end{pmatrix} - \begin{pmatrix} 25000 & 5 & 10 & \textcolor{blue}{5} & 0 & \frac{5}{3} & 0 \end{pmatrix} \\
        &= \begin{pmatrix} 1000 & -2 & -3 & \textbf{\textcolor{green!60!black}{0}} & 0 & -\frac{5}{3} & 1 \end{pmatrix}
    \end{align*}
\end{itemize}

\vspace{0.2cm}
\textbf{Tính toán dòng hệ số ước lượng $\Delta_j$ mới bằng ma trận:} \\
Tại Bảng 2, cơ sở mới là $\{x_4, x_3, x_6\}$. Vector hệ số hàm mục tiêu ứng với cơ sở này là $C_b = \begin{pmatrix} 0 \\ 15 \\ 0 \end{pmatrix}$. 
Hệ số ước lượng được tính theo công thức tổng quát $\Delta_j = C_b^T \times A_j - c_j$:
\begin{itemize}
    \item Giá trị hàm mục tiêu mới: $Z = C_b^T \times PA = 0(1000) + 15(5000) + 0(1000) = 75000$.
    \item $\Delta_1 = \begin{pmatrix} 0 & 15 & 0 \end{pmatrix} \begin{pmatrix} 1 \\ 1 \\ -2 \end{pmatrix} - 11 = 15 - 11 = 4$.
    \item $\Delta_2 = \begin{pmatrix} 0 & 15 & 0 \end{pmatrix} \begin{pmatrix} -4 \\ 2 \\ -3 \end{pmatrix} - 9 = 30 - 9 = 21$.
    \item $\Delta_5 = \begin{pmatrix} 0 & 15 & 0 \end{pmatrix} \begin{pmatrix} -4/3 \\ 1/3 \\ -5/3 \end{pmatrix} - 0 = 5$.
    \item Các biến cơ sở $x_3, x_4, x_6$ luôn có $\Delta_j = 0$.
\end{itemize}

% BẢNG ĐƠN HÌNH 2 (Bạn giữ nguyên đoạn code kẻ bảng ở đây)

\vspace{0.3cm}
\textbf{Giải thích chi tiết phần Kết luận:} \\
Theo tiêu chuẩn tối ưu của thuật toán đơn hình đối với bài toán tìm Max: nếu tất cả các hệ số ước lượng $\Delta_j \ge 0$ thì phương án cơ bản hiện hành là phương án tối ưu. 
Tại Bảng 2, ta có tập các hệ số ước lượng $\Delta = (4, 21, 0, 0, 5, 0)$ đều mang giá trị dương hoặc bằng $0$, do đó thỏa mãn hoàn toàn tiêu chuẩn dừng của thuật toán.

Từ cột Phương Án (PA) của Bảng 2, ta đọc được lời giải:
\begin{itemize}
    \item Các biến nằm trong cơ sở mang giá trị bằng đúng cột PA tương ứng: $x_3 = 5000, x_4 = 1000, x_6 = 1000$.
    \item Các biến không nằm trong cơ sở (vẫn ở ngoài) thì mang giá trị bằng $0$: $x_1 = 0, x_2 = 0, x_5 = 0$.
\end{itemize}
Chiếu theo bài toán kinh tế gốc, công ty không cần sản xuất sản phẩm loại A ($x_1=0$) và loại B ($x_2=0$), mà chỉ cần dồn toàn lực sản xuất 5000 sản phẩm loại C ($x_3=5000$). Khi đó, tiền lãi đạt mức cực đại là $75.000$ nghìn đồng (hay $75$ triệu đồng). Biến phụ $x_4 = 1000$ và $x_6 = 1000$ thể hiện rằng nguyên liệu I còn dư 1000g và nguyên liệu III còn dư 1000g.

\vspace{0.8cm}
\subsection{Lập bài toán đối ngẫu (D)}
\textbf{c. Lập bài toán đối ngẫu (D) của bài toán (P)}

Để lập bài toán đối ngẫu một cách chính xác mà không bị nhầm lẫn về dấu, ta sử dụng phương pháp Bảng cấu trúc ánh xạ. Dưới đây là quá trình chuyển đổi từ Bài toán gốc (Max) sang Bài toán đối ngẫu (Min):

\begin{table}[H]
    \centering
    \renewcommand{\arraystretch}{1.5}
    \resizebox{\textwidth}{!}{
    \begin{tabular}{|c|c c c c c c c c|c c c|c|}
        \hline
        \textbf{Hàm/RBC} & \multicolumn{8}{c|}{\textbf{Hệ số \& Biến số}} & \multicolumn{3}{c|}{\textbf{Dấu RBC}} & \textbf{Hệ số $b$} \\
        \hline
        $Z (\max)$ & \TL{c_1}{11} & $x_1$ & $+$ & \TL{c_2}{9} & $x_2$ & $+$ & \TL{c_3}{15} & $x_3$ & \multicolumn{3}{c|}{} & \\
        \hline
        $R_1$ & \TL{a_{11}}{5} & $x_1$ & $+$ & \TL{a_{12}}{4} & $x_2$ & $+$ & \TL{a_{13}}{4} & $x_3$ & $\ge$ & \khoanh{$\le$} & $=$ & \TL{b_1}{21000} \\
        \hline
        $R_2$ & \TL{a_{21}}{3} & $x_1$ & $+$ & \TL{a_{22}}{6} & $x_2$ & $+$ & \TL{a_{23}}{3} & $x_3$ & $\ge$ & \khoanh{$\le$} & $=$ & \TL{b_2}{15000} \\
        \hline
        $R_3$ & \TL{a_{31}}{3} & $x_1$ & $+$ & \TL{a_{32}}{7} & $x_2$ & $+$ & \TL{a_{33}}{5} & $x_3$ & $\ge$ & \khoanh{$\le$} & $=$ & \TL{b_3}{26000} \\
        \hline
        \multirow{3}{*}{\shortstack{Dấu \\ Biến số}} 
        & \multicolumn{2}{c}{$x_1$} & \multicolumn{6}{l|}{\khoanh{$\ge 0$} \quad $\le 0$ \quad Tùy ý} & \multicolumn{4}{c|}{} \\
        & \multicolumn{2}{c}{$x_2$} & \multicolumn{6}{l|}{\khoanh{$\ge 0$} \quad $\le 0$ \quad Tùy ý} & \multicolumn{4}{c|}{} \\
        & \multicolumn{2}{c}{$x_3$} & \multicolumn{6}{l|}{\khoanh{$\ge 0$} \quad $\le 0$ \quad Tùy ý} & \multicolumn{4}{c|}{} \\
        \hline
    \end{tabular}
    }
    \caption{Bảng phân rã cấu trúc bài toán Gốc (P)}
\end{table}

\vspace{0.3cm}
Từ Bảng phân rã Gốc, ta bốc các hệ số dựa vào ký hiệu để điền vào Bảng cấu trúc Đối ngẫu. Theo quy tắc ánh xạ, vì bài toán gốc là Max, các biến gốc $x_j \ge 0$ sẽ cho ra dấu $\ge$ ở RBC đối ngẫu; các RBC gốc dấu $\le$ sẽ cho ra biến đối ngẫu $y_i \ge 0$. Các dấu sai quy tắc sẽ bị gạch bỏ.

\begin{table}[H]
    \centering
    \renewcommand{\arraystretch}{1.5}
    \resizebox{\textwidth}{!}{
    \begin{tabular}{|c|c c c c c c c c|c c c|c|}
        \hline
        \textbf{Hàm/RBC} & \multicolumn{8}{c|}{\textbf{Hệ số \& Biến số}} & \multicolumn{3}{c|}{\textbf{Dấu RBC ĐN}} & \textbf{Hệ số $c$} \\
        \hline
        $Z' (\min)$ & \TL{b_1}{21000} & $y_1$ & $+$ & \TL{b_2}{15000} & $y_2$ & $+$ & \TL{b_3}{26000} & $y_3$ & \multicolumn{3}{c|}{} & \\
        \hline
        $R'_1$ & \TL{a_{11}}{5} & $y_1$ & $+$ & \TL{a_{21}}{3} & $y_2$ & $+$ & \TL{a_{31}}{3} & $y_3$ & \khoanh{$\ge$} & $\cancel{\le}$ & $\cancel{=}$ & \TL{c_1}{11} \\
        \hline
        $R'_2$ & \TL{a_{12}}{4} & $y_1$ & $+$ & \TL{a_{22}}{6} & $y_2$ & $+$ & \TL{a_{32}}{7} & $y_3$ & \khoanh{$\ge$} & $\cancel{\le}$ & $\cancel{=}$ & \TL{c_2}{9} \\
        \hline
        $R'_3$ & \TL{a_{13}}{4} & $y_1$ & $+$ & \TL{a_{23}}{3} & $y_2$ & $+$ & \TL{a_{33}}{5} & $y_3$ & \khoanh{$\ge$} & $\cancel{\le}$ & $\cancel{=}$ & \TL{c_3}{15} \\
        \hline
        \multirow{3}{*}{\shortstack{Dấu Biến \\ Đối ngẫu}} 
        & \multicolumn{2}{c}{$y_1$} & \multicolumn{6}{l|}{\khoanh{$\ge 0$} \quad $\cancel{\le 0}$ \quad $\cancel{\text{Tùy ý}}$} & \multicolumn{4}{c|}{} \\
        & \multicolumn{2}{c}{$y_2$} & \multicolumn{6}{l|}{\khoanh{$\ge 0$} \quad $\cancel{\le 0}$ \quad $\cancel{\text{Tùy ý}}$} & \multicolumn{4}{c|}{} \\
        & \multicolumn{2}{c}{$y_3$} & \multicolumn{6}{l|}{\khoanh{$\ge 0$} \quad $\cancel{\le 0}$ \quad $\cancel{\text{Tùy ý}}$} & \multicolumn{4}{c|}{} \\
        \hline
    \end{tabular}
    }
    \caption{Bảng cấu trúc bài toán Đối ngẫu (D) sau khi ánh xạ}
\end{table}

\textbf{Kết luận:} Dựa vào bảng ánh xạ trên, ta viết được mô hình của bài toán đối ngẫu (D) hoàn chỉnh như sau:
\begin{alignat*}{4}
    Z' = \quad & 21000y_1 & {}+{} & 15000y_2 & {}+{} & 26000y_3 & \rightarrow \min \\
    \text{s.t.} \quad & 5y_1 & {}+{} & 3y_2 & {}+{} & 3y_3 & \ge 11 \\
    & 4y_1 & {}+{} & 6y_2 & {}+{} & 7y_3 & \ge 9 \\
    & 4y_1 & {}+{} & 3y_2 & {}+{} & 5y_3 & \ge 15 \\
    & \rlap{$y_1, y_2, y_3 \ge 0$}
\end{alignat*}